\chapter{MCP and Agent Architecture}
\label{ch:mcp}

This chapter dives into the specific integration between the Model
Context Protocol (MCP) and the Groq-hosted LLM agent that drives the
platform. The material is organised bottom-up: protocol primitives,
tool discovery, tool invocation, agent decision flow, and
extensibility.

\section{MCP in One Page}
\label{sec:mcp-onepager}

MCP \citep{mcpSpec2024, anthropic2024mcp} is a JSON-RPC 2.0 protocol
in which an \emph{MCP host} (typically an LLM application) connects
to one or more \emph{MCP servers}, each of which exposes some
combination of tools, resources, and prompts. Communication happens
through a \emph{transport}. The reference implementation uses stdio
or WebSockets, but FastMCP \citep{fastmcp} also provides
\code{FastMCPTransport}, which runs the entire client/server exchange
as ordinary \code{async} method calls inside the same Python
process. \projname{} uses \code{FastMCPTransport} exclusively.

\section{The In-Process Pool}
\label{sec:in-process-pool}

Figure~\ref{fig:pool} shows the \code{MCPClientPool} at startup. The
pool wraps each FastMCP server in a \code{Client(FastMCPTransport(server))},
enters the client's async context in an \code{AsyncExitStack}, and
records the state.

\begin{figure}[H]
\centering
\begin{tikzpicture}[node distance=0.55cm and 1.2cm, font=\small\sffamily]
  \node[al-node, minimum width=3cm]                     (a) {App starts};
  \node[al-node, right=of a, minimum width=3.5cm]       (b) {Instantiate \\ \code{MCPClientPool}};
  \node[al-node, right=of b, minimum width=3.5cm]       (c) {\code{register\_default\_} \\ \code{servers()}};
  \node[al-node, below=of c, minimum width=3.5cm]       (d) {For each server: \\ open Client};
  \node[al-node, left=of d, minimum width=3.5cm]        (e) {Rebuild tool + \\ prompt registries};
  \node[al-node, left=of e, minimum width=3cm]          (f) {Pool ready};

  \draw[al-flow] (a) -- (b);
  \draw[al-flow] (b) -- (c);
  \draw[al-flow] (c) -- (d);
  \draw[al-flow] (d) -- (e);
  \draw[al-flow] (e) -- (f);
\end{tikzpicture}
\caption{Pool startup sequence.}
\label{fig:pool}
\end{figure}

Once startup finishes, the pool exposes four public interfaces:

\begin{itemize}[leftmargin=1.4em]
  \item \code{tool\_schemas} --- a list of OpenAI-compatible function
        schemas ready to hand to the LLM.
  \item \code{call\_tool(name, arguments, progress\_handler)} ---
        routes to the correct client and forwards progress events.
  \item \code{get\_prompt(name, arguments)} --- retrieves a
        parameterised prompt template.
  \item \code{snapshot()} --- an inspectable dictionary of service
        state, tool counts, and prompt metadata (used by the sidebar
        health panel).
\end{itemize}

\section{Tool Discovery}
\label{sec:discovery}

Immediately after all clients connect, \code{\_rebuild\_registries}
loops through the clients and calls \code{list\_tools},
\code{list\_resources} (best-effort), and \code{list\_prompts}
(best-effort). For each tool it constructs an OpenAI-compatible
schema of the form
\begin{lstlisting}[caption={Schema shape passed to the Groq API.},
                    label={lst:openai-schema}]
{
  "type": "function",
  "function": {
    "name":        "<tool_name>",
    "description": "<tool.description>",
    "parameters":  <tool.inputSchema>       # JSON Schema
  }
}
\end{lstlisting}
The pool preserves the natural-language tool description generated by
FastMCP from the Python docstring; this is what the LLM sees when it
decides whether a tool is relevant.

\section{Tool Invocation Sequence}
\label{sec:invocation}

Figure~\ref{fig:seq} shows a single tool call from the LLM's
perspective. The sequence is compressed --- in a real turn each of
the arrows is streamed piecewise.

\begin{figure}[H]
\centering
\begin{tikzpicture}[node distance=0.4cm and 1.4cm, font=\small\sffamily]
  \node[al-node, minimum width=2.2cm] (ui) {UI};
  \node[al-node, right=of ui, minimum width=2.8cm] (or) {Orchestrator};
  \node[al-node, right=of or, minimum width=2cm] (llm) {LLM};
  \node[al-node, right=of llm, minimum width=2.5cm] (po) {Pool};
  \node[al-node, right=of po, minimum width=2.5cm] (srv) {MCP tool};

  \node[below=1.6cm of ui, minimum width=0.1cm] (ui2){};
  \node[below=1.6cm of or, minimum width=0.1cm] (or2){};
  \node[below=1.6cm of llm, minimum width=0.1cm] (llm2){};
  \node[below=1.6cm of po, minimum width=0.1cm] (po2){};
  \node[below=1.6cm of srv, minimum width=0.1cm] (srv2){};

  \draw[al-flow] (ui.south)  -- ($(or.south)+(0,-0.2)$) node[midway,above,font=\tiny]{user query};
  \draw[al-flow] (or.south)  -- ($(llm.south)+(0,-0.4)$) node[midway,above,font=\tiny]{messages + tool schemas};
  \draw[al-flow, dashed] (llm.south) -- ($(or.south)+(0,-0.7)$) node[midway,above,font=\tiny]{tool\_call: name+args};
  \draw[al-flow] (or.south)  -- ($(po.south)+(0,-1.0)$) node[midway,above,font=\tiny]{call\_tool(name, args)};
  \draw[al-flow] (po.south)  -- ($(srv.south)+(0,-1.3)$) node[midway,above,font=\tiny]{routed invocation};
  \draw[al-flow, dashed] (srv.south) -- ($(po.south)+(0,-1.5)$) node[midway,above,font=\tiny]{result};
  \draw[al-flow, dashed] (po.south) -- ($(or.south)+(0,-1.5)$);
  \draw[al-flow] (or.south)  -- ($(llm.south)+(0,-1.5)$) node[midway,above,font=\tiny]{tool result appended};
  \draw[al-flow, dashed] (or.south) -- ($(ui.south)+(0,-1.5)$) node[midway,above,font=\tiny]{tool\_end event};

\end{tikzpicture}
\caption{Single tool-call round trip. Dashed arrows indicate
streamed responses.}
\label{fig:seq}
\end{figure}

\section{Progress Bridging}
\label{sec:progress}

MCP tools can report progress by calling
\code{ctx.report\_progress(current, total, message)}. FastMCP
delivers these to a progress handler supplied by the client. The
orchestrator installs a handler that pushes each event into an
\code{asyncio.Queue}. During tool execution, the orchestrator polls
the queue with a short timeout and yields any accumulated events as
\code{tool\_progress} events. This is what enables the ``training
model 3/4 --- 42\% complete'' line that the Streamlit UI renders in
place.

\section{The Agent Decision Flow}
\label{sec:agent-flow}

The agent is a single-turn planner in the ReAct style
\citep{yao2023react}: it receives (i)~a system prompt that describes
its role and the tool catalog, (ii)~the conversation history, and
(iii)~the current user query, and it emits either plain text or a
tool call. The system prompt (constructed by
\code{orchestrator/loop.py::\_build\_messages}) instructs the model
to:

\begin{itemize}[leftmargin=1.4em]
  \item Prefer the \emph{active dataset} named in context when the
        user says ``my data'' or ``the dataset''.
  \item Never invent numerical claims --- always ground them in tool
        results.
  \item After a bake-off succeeds, proactively compute SHAP values
        with the returned \code{model\_id} and
        \code{x\_train\_sample\_id}.
  \item Prefer short paragraphs and light bullets.
\end{itemize}

Multi-step behaviour is achieved by allowing the LLM to emit
\emph{multiple} tool calls per turn and by permitting up to
\code{MAX\_TOOL\_ITERATIONS} loops, so a request such as ``run EDA
and then train a model'' typically fires 3--5 tools in a single
user turn.

\section{Security Considerations}
\label{sec:security-mcp}

Because \projname{} is a single-user local tool, the threat model is
lightweight, but a few observations are worth recording.

\begin{itemize}[leftmargin=1.4em]
  \item \textbf{Prompt injection.} The tool descriptions are trusted
        (they are versioned code), but the \emph{results} of tools
        are attacker-influenced (they may contain data from an
        uploaded CSV). The system does not currently sanitise
        tool results before feeding them back to the LLM; a
        malicious CSV cell could in principle contain instructions
        aimed at the model. In practice this is mitigated by
        showing the entire tool result to the user in an expander;
        no unattended automation occurs.
  \item \textbf{Filesystem writes.} Every write goes through
        \code{workspace.save\_artifact}, which sanitises the filename
        via \code{\_safe\_stem} (only alphanumerics, dashes,
        underscores, and dots survive). This prevents path traversal
        in \code{data/artifacts/}.
  \item \textbf{Arbitrary code execution.} The platform never
        \code{exec}s LLM-produced code. The LLM can only invoke
        pre-defined tools with typed arguments.
  \item \textbf{Network.} The only outbound network call is to Groq;
        no inbound network is opened.
\end{itemize}

\section{Extensibility}
\label{sec:mcp-extensibility}

The MCP layer is trivially extensible:

\begin{enumerate}[leftmargin=1.4em]
  \item Add a new tool by decorating a function with
        \code{@mcp.tool} in one of the existing servers.
  \item Add a new server by creating a new
        \file{mcp\_servers/<name>.py} module and registering its
        \code{mcp} instance in
        \code{register\_default\_servers()}.
  \item Add a prompt template by decorating a function with
        \code{@mcp.prompt}; it becomes visible under
        \code{pool.prompts\_meta}.
\end{enumerate}

No changes to \file{orchestrator/loop.py} or
\file{app/streamlit\_app.py} are required. The new tool will appear
in the LLM's tool schema list on the next start-up.

\section{Comparison to HTTP-Based MCP}
\label{sec:vs-http}

MCP is more commonly deployed over stdio or Streamable HTTP. The
in-process transport has three practical advantages for a research
platform:

\begin{itemize}[leftmargin=1.4em]
  \item No port management, no shutdown races, no partial startup.
  \item Zero-copy Python object passing between client and server.
  \item Perfect traceability in a Python debugger --- each tool
        invocation is a normal call stack.
\end{itemize}

The disadvantages (no cross-language servers, no independent scaling
per service) are acceptable for a single-user tool. Because the
protocol contract is unchanged, the servers can be re-hosted over
HTTP simply by swapping the transport in \code{register\_default\_servers}.
